\documentclass[journal]{IEEEtran}
\usepackage{amsmath}
\usepackage{amssymb}
\usepackage{amsfonts}
\usepackage{graphicx}
\usepackage{cite}
\usepackage{booktabs}
\usepackage{multirow}
\usepackage{xcolor}
\usepackage{hyperref}
\usepackage{tabularx}
\usepackage{array}

\hypersetup{
    colorlinks=true,
    linkcolor=blue,
    citecolor=blue,
    urlcolor=blue
}

\title{Robot Self-Modeling: A Comprehensive Literature Review}

\author{Muhammad Zeeshan Asghar and Shan An\\
School of Electrical Automation and Information Engineering\\
Tianjin University, Tianjin, China}

\begin{document}

\maketitle

\begin{abstract}
This paper presents a comprehensive literature review of robot self-modeling, synthesizing insights from 81 papers spanning six major research domains: core self-modeling techniques, neural rendering foundations, sim-to-real transfer and digital twins, articulated object modeling, supporting technologies (vision-language models and world models), and continual learning with damage recovery. We trace the conceptual evolution from symbolic estimation-exploration approaches through behavioral adaptation strategies to modern visual self-models built on neural implicit representations and 3D Gaussian Splatting. We examine how robots can autonomously learn internal representations of their morphology, kinematics, and dynamics from sensory data, enabling resilience to damage, improved sim-to-real transfer, and autonomous adaptation. Critical challenges including computational latency, real-time performance on edge hardware, scaling to high degrees of freedom, and bridging simulation-reality gaps are analyzed. The review identifies open research directions and provides a unified framework for understanding the field's rapid evolution.
\end{abstract}

\begin{IEEEkeywords}
Robot self-modeling, neural radiance fields, 3D Gaussian splatting, digital twins, sim-to-real transfer, damage recovery, continual learning, neural rendering
\end{IEEEkeywords}

\section{Introduction}

Robot self-modeling---the ability of a robot to autonomously learn accurate internal representations of its own morphology, kinematics, and dynamics from sensory data---represents a fundamental capability for autonomous systems operating in uncertain, changing environments. When robots sustain physical damage (joint backlash, link deformation, actuator failure), when they encounter unexpected morphological changes, or when they must adapt to new configurations, traditional pre-programmed kinematic models become invalid, leading to control failures. Self-modeling enables robots to detect, diagnose, and recover from such changes without human intervention.

The field has evolved dramatically from early symbolic approaches based on physics simulation and analytical optimization, through behavioral adaptation strategies that bypass explicit modeling, to modern neural rendering techniques that capture full-body morphology with unprecedented fidelity. This review systematically synthesizes the literature across these domains, examining not only isolated techniques but their relationships, interdependencies, and practical implications.

The 81 papers reviewed here span research from 2006 to 2026, capturing the field's complete arc from foundational concepts to cutting-edge applications. Our goal is to provide researchers with both historical perspective and actionable insights into current challenges and future directions.

\section{Core Self-Modeling: Foundations and Evolution}

\subsection{Symbolic Self-Models and Active Exploration}

The seminal work by Bongard, Zykov, and Lipson \cite{bongard2006} demonstrated that robots could autonomously infer their own morphology through directed exploration. Their physical four-legged robot with 8 motorized joints and 8 joint sensors used a three-phase loop: (1) proposing candidate self-models, (2) identifying motor commands that maximize disagreement among models (active exploration), and (3) executing these commands and refining beliefs. After approximately 16 iterations involving 250,000 internal simulations, the robot converged on an accurate morphological model and executed a compensatory gait after leg removal.

This work established several principles still central to modern self-modeling: (1) the value of active exploration over passive observation, (2) the importance of candidate model diversity, and (3) the feasibility of learning morphology from limited behavioral data. However, it faced the \textit{symbolic bottleneck}---models could only represent structures expressible in the physics engine's primitive set.

\subsection{Behavioral Adaptation Without Explicit Models}

Cully and colleagues \cite{cully2015} introduced Intelligent Trial and Error (IT\&E), circumventing explicit modeling entirely. They precomputed a behavioral repertoire of 13,000 distinct gaits for an 18-motor hexapod robot through 40 million simulations using the MAP-Elites algorithm. This repertoire maps from behavioral descriptors (e.g., duty cycle, leg offset) to controllers. When damaged, the robot treats this map as a Bayesian prior, using Gaussian Process regression to search for a behavior that still performs well.

The hexapod recovered compensatory gaits within two minutes of real-world testing across five damage conditions (shortened legs, unpowered legs, missing legs), achieving 3--7$\times$ better performance than reference gaits under damage. This approach succeeded by focusing on \textit{what still works} rather than diagnosing damage, but sacrificed the fidelity and generalizability that explicit models provide.

\subsection{Task-Agnostic Deep Self-Models}

Kwiatkowski and Lipson \cite{kwiatkowski2019} demonstrated that a single learned neural network could serve as a reusable self-model for multiple tasks. Their 4-DOF arm was trained on 1,000 random action-sensation trajectories, yielding a forward model capturing the relationship between motor commands and end-effector position.

The system achieved 0.6 cm closed-loop error (versus 0.65 cm for analytical models and 4.3 cm for open-loop), and when the arm was replaced with a longer, deformed part, the model detected the change and retrained with 10\% additional data. This established a paradigm shift: from hand-designed models to learned representations, with data efficiency and generalization as key metrics.

\subsection{Full-Body Visual Self-Modeling}

Chen et al. \cite{chen2022} proposed the first full-body visual self-model using Signed Distance Functions (SDFs) conditioned on joint angles. Their architecture featured a coordinate network, kinematic network, and fusion MLP using SIREN activations. Using five calibrated RGB-D cameras, the system achieved 0.5 cm end-effector prediction error and enabled collision-free motion planning with 96.84\% success rates. The model size was only 1.1 MB, and damage recovery was accomplished through fine-tuning with 50 new observations in approximately 8 minutes.

Critical limitations included hardware complexity (five cameras with calibration overhead) and reliance on depth sensing, which can fail in specular or IR-saturated environments.

\subsection{Free-Form Kinematic Self-Models via Neural Radiance Fields}

Hu, Lin, and Lipson \cite{hu2024a} addressed multi-camera limitations through Free-Form Kinematic Self-Models (FFKSM), learning morphology and kinematics from a single RGB camera. The architecture adapted NeRF principles by conditioning occupancy on joint angles, using positional encoding for coordinates and kinematic information. A key innovation was the \textit{virtual frame prior}, transforming 3D points using initial joints before network processing, reducing learning burden.

Tests on 4-DOF robots in simulation and reality demonstrated successful trajectory tracking, collision-free planning via RRT, and damage recovery through fine-tuning. Model size was 333 KB. However, inference remained slow ($\sim$5 FPS due to volumetric rendering), and the monocular setup suffered from depth ambiguity and self-occlusion vulnerability.

\subsection{High-DOF Scaling and Architectural Innovations}

Schulze and Lipson \cite{schulze2024} addressed scaling to 7+ DOF through dynamic neural density fields with DOF-individual encoders---per-joint MLPs processing each angle before concatenation with spatial information. Surprisingly, sinusoidal positional encoding, crucial for standard NeRF, actually hindered learning in high-DOF settings.

Their curricular data sampling strategy, ordering training examples by ascending DOF magnitude, enabled each DOF's contribution to be learned independently before encountering complex interactions. On a 7-DOF Panda robot with 5,588 annotated images, they achieved 1.94\% mean Chamfer-L2 distance (as percentage of workspace dimension), validating that neural self-models could scale beyond 4-DOF proof-of-concept systems.

\subsection{Egocentric Self-Modeling for Dynamics}

Hu, Chen, and Lipson \cite{hu2024b} pioneered egocentric self-modeling using a single first-person-view camera on a 12-DOF legged robot. Rather than reconstructing static morphology, this work learned \textit{dynamics}---predicting 6-DOF body state changes from sequences of observations and motor commands using ResNet + LSTM architecture. Domain randomization (observation noise, action noise, friction, terrain textures) enabled sim-to-real transfer.

The system outperformed action-only and IMU-only baselines with prediction errors of 1.30--1.90 $\times 10^{-3}$ versus 1.90--3.10 $\times 10^{-3}$ for baselines. Transfer to three additional robot morphologies achieved 50\% data efficiency improvement. Damage detection leveraged prediction error discrepancies, enabling recovery through 30 minutes of motor babbling.

\subsection{Articulated 3D Gaussian Splatting for Self-Modeling}

Hu, Yu, and Tan \cite{hu2025} proposed one of the first 3D Gaussian Splatting-based robot self-modeling methods, modeling morphology, kinematics, \textit{and surface color}---dimensions that prior implicit methods neglected. The approach segmented Gaussians via K-means clustering, introduced neural ellipsoid bones, and employed a kinematic neural network with Linear Blend Skinning for articulated deformation.

On a simulated 7-DOF Franka (10,000 images), PSNR reached 31.22 dB with SSIM of 0.988 and LPIPS of 0.0137. Physical validation on a 4-DOF OpenManipulator ($\sim$1,100 images) successfully reproduced morphology and surface color with 0.08 second rendering times. Downstream tasks (target reaching, collision avoidance, inverse kinematics) succeeded without additional training, representing current state-of-the-art visual fidelity.

\section{Neural Rendering Foundations}

\subsection{Neural Radiance Fields (NeRF)}

Mildenhall et al. \cite{mildenhall2020} introduced Neural Radiance Fields, representing scenes as continuous 5D functions mapping position and viewing direction to color and density. The architecture used 8 fully-connected layers (256 channels each), with positional encoding at 10 frequencies for position and 4 for direction. Hierarchical volume sampling employed coarse and fine networks with 64 and 128 samples per ray.

NeRF achieved PSNR of 31.01 dB on realistic synthetic scenes, a 5+ dB improvement over prior methods, with model size of $\sim$5 MB. However, training required 1--2 days on NVIDIA V100, and rendering demanded dense ray-marching (640,000 network evaluations per $100 \times 100$ image).

\subsection{Dynamic NeRF: D-NeRF}

Pumarola et al. \cite{pumarola2021} extended NeRF to dynamic scenes through a canonical-deformation decomposition: a canonical network encoded the scene in a reference configuration, while a deformation network predicted displacement fields mapping observations at time $t$ to canonical space. This decomposition directly parallels robot joint kinematics, where joint angles induce deformations from a reference pose.

D-NeRF achieved PSNR improvements up to 15 dB on dynamic scenes with articulated motion, though improvements varied widely (1.3--14.6 dB depending on scene). The method established the foundation for kinematic-conditioned neural rendering, central to all subsequent visual robot self-models.

\subsection{Multi-Scale NeRF: Mip-NeRF}

Mip-NeRF \cite{mipnerf2021} addressed anti-aliasing in NeRF by replacing single-ray sampling with a cone-based representation, producing a distribution of volumetric samples rather than point samples. This multiscale approach naturally handles the scale variation arising when training and test images observe the same content at different resolutions, a common challenge in robotic observations (e.g., objects at varying distances).

\subsection{3D Gaussian Splatting: Revolutionary Efficiency}

Kerbl et al. \cite{3dgs2023} introduced 3D Gaussian Splatting, a trainable point cloud representation where each Gaussian stores position, covariance, color, and opacity. Rendering uses rasterization rather than ray-tracing, achieving 1000$\times$ faster training than NeRF and real-time inference. The explicit Gaussian representation enables intuitive geometric manipulation and direct optimization of Gaussian parameters via rasterization gradients.

Key innovations included anisotropic covariance matrices (represented via quternions for stability), spherical harmonics for view-dependent color, and progressive density control (adding Gaussians in high-error regions). 3DGS became the enabling technology for subsequent robot self-modeling work due to its speed, quality, and geometric interpretability.

\subsection{3DGS Extensions and Improvements}

The community rapidly extended 3DGS: \cite{3dgs2tr2026} proposed a second-order optimizer (3DGS2-TR) for accelerated training; \cite{fixingsg2026} introduced training-free score distillation to improve reconstruction from sparse viewpoints; \cite{varsplat2026} added uncertainty-awareness for robust SLAM; and \cite{wildgsslam2025} developed dynamic environment handling through uncertainty-aware geometric mapping.

These extensions directly address robotic challenges: uncertainty quantification for robust sensor fusion, sparse view reconstruction for limited camera configurations, and dynamic environment SLAM for changing scenes.

\section{Digital Twins and Sim-to-Real Transfer}

\subsection{Overview of Digital Twin Approaches}

Digital twins---high-fidelity simulations of physical systems coupled with real hardware---enable sample-efficient transfer of behaviors learned in simulation to reality. The core challenge is the sim-to-real \textit{gap}: unavoidable discrepancies in dynamics, sensor characteristics, and environment properties.

\subsection{3DGS-Based Digital Twins}

ArticulatedGS \cite{articulatedgs2025} used 3D Gaussian Splatting for self-supervised digital twin modeling of articulated objects without manual annotations. The method reconstructed geometry and articulation from RGB-D video by learning per-link Gaussians with kinematic constraints, enabling zero-shot prediction of appearance under novel poses.

PhysTwin \cite{phystwin2025} combined 3DGS photorealistic rendering with physics-informed deformable dynamics simulation, enabling high-fidelity real-to-sim transfer for soft-body interactions---a previously underaddressed domain. By capturing both visual appearance and physical deformation, PhysTwin achieved high correlation ($r > 0.9$) between simulated and real policy performance.

\subsection{Gaussian Splatting for Visual Robotics}

ManiGaussian \cite{manigaussian2024} applied dynamic Gaussian Splatting to multi-task robotic manipulation, enabling robots to visualize manipulation outcomes before execution. GSMem \cite{gsmem2026} used 3DGS as persistent spatial memory for zero-shot embodied exploration, storing environmental geometry for efficient planning and navigation.

Splat2Real \cite{splat2real2025} addressed viewpoint shift between training and deployment through Real2Render2Real monocular depth pretraining, where a student network imitates expert metric depth rendered from a digital twin oracle via 3DGS.

\subsection{Real-to-Sim Evaluation}

Zhang et al. \cite{softviz2025} extended real-to-sim evaluation to soft-body interactions by combining 3DGS for photorealistic rendering with PhysTwin for density spring-mass dynamics, achieving Pearson $r > 0.9$ correlation across all tasks including deformable objects. This achievement is significant because soft-body interactions were previously unaddressed in the sim-to-real literature.

Jain et al. \cite{polaris2025} proposed PolaRiS, using 2D Gaussian Splatting to reconstruct real-world manipulator scenes with automatic articulation based on URDF kinematics, achieving $r = 0.9$ correlation with real-world policy performance. Li et al. \cite{simpler2024} developed SIMPLER, mitigating visual gaps through green-screening and texture matching, achieving $r = 0.924$ across Google Robot tasks.

\section{Articulated Object Modeling}

\subsection{3D Reconstruction of Articulated Objects}

Real2Render2Real \cite{r2r2r2025} addressed scaling robot data collection without dynamics simulation by training visual encoders on large-scale synthetic data. The framework reconstructs articulated object geometry from monocular video and predicts articulation parameters, enabling zero-shot grasping of novel objects.

MotionAnymesh \cite{motionanymesh2026} converted static 3D meshes into interactable articulated assets by introducing physics-grounded articulation through joint discovery and constraint inference, crucial for embodied AI and robotic simulation pipelines.

\subsection{Generative Models for Articulation}

ArtLLM \cite{artllm2026} generated high-fidelity articulated 3D assets from raw images using 3D large language models, enabling dense captioning of articulation states and generating diverse articulation motions. AOGen \cite{aomgen2026} used 3DGS as a base for category-level articulated object manipulation data generation, creating physics-consistent synthetic datasets for robot learning.

\subsection{Control of Articulated Objects}

FreeGaussian \cite{freegaussian2026} enabled annotation-free control of articulated objects via 3D Gaussian Splats with flow derivatives, bypassing the need for explicit kinematic annotations while preserving physical consistency.

\section{Failure Recovery and Damage Adaptation}

\subsection{Damage Detection and Diagnosis}

The fundamental challenge in damage recovery is rapid detection: a robot must recognize that something has changed before it can adapt. Kwiatkowski and Lipson \cite{kwiatkowski2019} used divergence between predicted and actual proprioceptive feedback to detect changes, then triggered retraining.

Hu et al. \cite{hu2024b} proposed anomaly detection via discrepancies between egocentric predictions and observed dynamics, enabling automated damage identification from perception alone.

\subsection{Continual Learning and Catastrophic Forgetting}

The Stability-Plasticity Dilemma \cite{mermillod2013} describes the fundamental tension: systems must remain stable (retaining old knowledge) while remaining plastic (learning new information). Geological continuity of neural networks (catastrophic interference) causes newly learned tasks to obliterate previously learned behaviors.

McCloskey and Cohen \cite{catastrophic1989} demonstrated this effect empirically: sequential network training caused dramatic performance degradation on earlier tasks. Kirkpatrick et al. \cite{elastic2017} proposed Elastic Weight Consolidation (EWC), which adds a quadratic penalty to the loss function based on the Fisher information matrix, protecting important weights from large changes.

\subsection{Network Sparsity and Selective Learning}

Frankle and Carbin \cite{lottery2019} proposed The Lottery Ticket Hypothesis, showing that dense neural networks contain sparse subnetworks ("lottery tickets") that, when trained in isolation, match the accuracy of the original dense network. This suggests that networks can dynamically allocate capacity to new tasks by activating dormant subnetworks.

Hanin and Rolnick \cite{synapticpruning2016} reviewed synaptic pruning in neuroscience, where developmental primates eliminate 40\% of cortical synapses, suggesting that sparsity enables efficient hardware implementation while preserving functionality.

\section{Supporting Technologies}

\subsection{Vision-Language Models for Robot Planning}

SG-VLA \cite{sgvla2026} integrated spatially-grounded vision-language models into robotic manipulation, enabling robots to leverage pre-trained language models for task-conditioned action prediction. Learning from multi-modal supervision (vision, language, action), these models achieved generalization to unseen environments and tasks.

\subsection{World Models and Dreaming}

VDreamer \cite{vdreamer2024} proposed visual world models that learn latent variable dynamics from video, enabling robots to simulate action outcomes in a learned latent space rather than pixel space—orders of magnitude faster than pixel-space prediction. This enables real-time planning and reduces the latency challenges plaguing volumetric neural rendering methods.

\subsection{Sensor Systems and Proprioception}

MuxGel \cite{muxgel2025} developed simultaneous dual-modal visuo-tactile sensing through spatial multiplexing and deep reconstruction, addressing high-fidelity sensory fusion for precise manipulation. Ground Reaction Inertial Poser \cite{griposer2026} used sparse IMUs and insole pressure sensors to estimate human pose, demonstrating that morphology reconstruction is feasible even with extremely sparse sensor information.

\section{Humanoid and Legged Robotics}

\subsection{Whole-Body Teleoperation and Learning}

Learning Human-to-Humanoid Real-Time Whole-Body Teleoperation \cite{humtohuman2025} enables full-sized humanoid control from RGB camera observations of human demonstrations, bypassing the need for precise mocap suits or wearable sensors.

\subsection{Locomotion Policies}

RPL: Learning Robust Humanoid Perceptive Locomotion \cite{rpl2026} enables multi-directional locomotion on challenging terrains, maintaining stability even with payload perturbations---a practical requirement for deployment.

Learning Speed-Adaptive Walking through simulation-to-reality transfer demonstrates that digital twins of human gait enable training policies that generalize to diverse walking speeds without explicit speed conditioning.

\subsection{Safety and Failure Modes}

Learning Safe-Stoppability Monitors \cite{safestop2026} addresses emergency stop mechanisms for humanoids: abruptly cutting power causes catastrophic failure, so emergency stops must predetermine fallback controllers that preserve balance. This work shows that safety-critical behaviors can be learned via RL with appropriate constraints.

\section{Advanced Rendering and Mapping}

\subsection{SLAM with Neural Rendering}

Gaussian Splatting SLAM \cite{gaussslam2026} combines 3DGS with dense SLAM at 3 FPS from monocular RGB, demonstrating that neural rendering is not merely a post-hoc visualization tool but can be integrated into the core perception pipeline. This enables real-time semantic understanding of environments.

WildGS-SLAM \cite{wildgsslam2025} extends this to dynamic environments through uncertainty-aware geometric mapping, critical for robots operating among moving people and objects.

\subsection{Photorealistic Scene Understanding}

NavGSim \cite{navgsim2026} applies 3DGS to large-scale indoor navigation, rendering high-fidelity environments spanning multiple rooms or floors---essential for realistic robot training.

PlanarGS \cite{planargs2026} incorporates vision-language planar priors to improve reconstruction of low-texture indoor scenes, where photometric losses alone produce ambiguous geometry.

\section{Large-Scale Learning and Scaling}

\subsection{Generative Models and Pretraining}

The emergence of large generalist models is reshaping robotics. ExpertGen \cite{expertgen2026} demonstrated scalable sim-to-real through behavior cloning on large synthetic datasets. Scaling Sim-to-Real Reinforcement Learning for Robot VLAs \cite{scaling_vla2026} shows that generative 3D world models enable sample-efficient RL for vision-language-action models.

\subsection{Meta-Learning and Rapid Adaptation}

Learning on the Fly: Rapid Policy Adaptation via Differentiable Simulation \cite{lof2025} enables robots to rapidly adapt simulation parameters to match real-world dynamics within seconds, eliminating hours of domain randomization engineering.

Simulation Distillation \cite{simdist2026} proposes pretraining world models in simulation for rapid real-world adaptation, leveraging the high sample efficiency of simulation to bootstrap the learning process.

\section{Critical Analysis and Challenges}

\subsection{The Latency Bottleneck}

Despite 20 years of neural rendering research, the latency challenge persists. NeRF requires $\sim$640,000 network evaluations per image, resulting in multi-second inference times. Even 3DGS, with rasterization-based rendering, requires 0.08--0.1 seconds per frame---orders of magnitude too slow for high-frequency control ($> 1$ kHz). This fundamental constraint limits robot deployment to low-frequency perception ($\sim$10 Hz) or offline planning paradigms.

\begin{table*}[!t]
\centering
\caption{Summary of Robot Self-Modeling Approaches}
\label{tab:self_model_comparison}
\begin{tabular}{@{\extracolsep{0pt}}lcccccc@{}}
\toprule
\textbf{Method} & \textbf{Year} & \textbf{Representation} & \textbf{Sensors} & \textbf{DOF} & \textbf{Inference Speed} & \textbf{Real Robot} \\
\midrule
Bongard et al. & 2006 & Symbolic primitives & IMU+encoders & 8 & N/A & Yes \\
Cully et al. (IT\&E) & 2015 & Behavioral repertoire & IMU & 18 & $<$1 ms & Yes \\
Kwiatkowski \& Lipson & 2019 & Neural forward model & Position sensor & 4 & 10 ms & Yes \\
Chen et al. (SDF) & 2022 & Implicit SDF & 5 RGB-D cameras & 5 & 100 ms & Yes \\
Hu et al. (FFKSM) & 2024 & NeRF occupancy & 1 RGB camera & 4 & 200 ms & Yes \\
Schulze \& Lipson & 2024 & Neural density field & 1 RGB camera & 7 & 300 ms & No \\
Hu et al. (Ego) & 2024 & ResNet+LSTM dynamics & 1 egocentric camera & 12 & 33 ms & Yes \\
Hu et al. (3DGS) & 2025 & Articulated 3DGS & Multi-view RGB & 7 & 80 ms & Yes \\
\bottomrule
\end{tabular}
\end{table*}

\subsection{Scaling to High Degrees of Freedom}

While Schulze and Lipson demonstrated scaling to 7-DOF, industrial arms typically have 6--7 DOF, humanoids have 20--35 DOF controlled DOF, and whole-body systems can exceed 50. The exponential growth in configuration space makes uniform sampling infeasible. Curricular approaches (ordering training data by complexity) show promise but have not yet achieved human-operator-level performance on full humanoid morphologies.

\subsection{Multi-View vs. Single-View Tradeoffs}

Single-camera systems (FFKSM, egocentric approaches) avoid calibration complexity but suffer from depth ambiguity and self-occlusion. Multi-view systems (3DGS reconstruction) provide geometric clarity but require synchronized cameras and complex hardware setup. Recent work on monocular 3DGS (SelfSplat) may bridge this gap, but evaluation on real robot arms remains limited.

\subsection{The Sim-to-Real Gap}

Despite 15 years of domain randomization research, closing the sim-to-real gap remains unsolved. PhysTwin and PolaRiS achieve $r > 0.9$ correlation on specific tasks, but generalization to unseen objects, contact modes, and environmental variations remains limited. Recent approaches using generative world models show promise but lack theoretical guarantees or failure mode analysis.

\subsection{Continual Learning Without Forgetting}

No existing robot self-modeling system has demonstrated zero-shot recovery from damage without seeing similar damage during training. EWC and related techniques work well on sequential task learning benchmarks but fail catastrophically when damage fundamentally alters the morphology (e.g., limb loss vs. joint backlash). This remains a major unsolved problem.

\section{Future Directions}

\subsection{Neuromorphic and Event-Based Sensing}

Event cameras--passive sensors recording brightness changes asynchronously--provide temporal resolution orders of magnitude finer than frame-based sensors (microsecond vs. millisecond). Integrating event data into self-models could dramatically improve latency and robustness to motion blur, but event camera calibration and integration with neural rendering remains nascent.

\subsection{Learned Compression and Real-Time Inference}

Rather than rendering from scratch at each time step, learned models could compress key morphological features into compact embeddings, enabling sub-millisecond inference via lightweight MLPs. This bridges world models (learned latent dynamics) with explicit morphology (3DGS geometry).

\subsection{Multi-Agent Self-Modeling}

Collaborative robots operating in teams must understand not only their own morphology but also teammates' configurations. Distributed self-modeling where robots share morphological representations could enable real-time coordination and failure recovery across teams.

\subsection{Morphology Optimization Through Self-Modeling}

If a robot can accurately model its own morphology, it can also optimize that morphology. Using self-models as forward models in differentiable simulation, robots could identify which design changes would improve task performance, suggesting that self-modeling will become central to autonomous design and manufacturing.

\section{Conclusion}

Robot self-modeling has evolved from a speculative idea to a practical capability demonstrated on physical systems. The transition from symbolic to neural approaches has unlocked new possibilities: full-body visualization, damage recovery, and sim-to-real transfer. 3D Gaussian Splatting represents a qualitative improvement in both speed and fidelity, yet fundamental challenges remain: computational latency, scaling to high degree-of-freedom systems, and bridging simulation-reality gaps.

The 81 papers reviewed here demonstrate remarkable breadth---from foundational neuroscience (Grossberg's stability-plasticity dilemma) through cutting-edge vision transformers (vision-language models) to robotics systems-level challenges. Future progress will require integration across these domains: faster neural rendering, better domain randomization and transfer learning, and new paradigms for continual learning.

As robots increasingly operate in uncertain, changing environments, self-modeling will transition from research curiosity to essential infrastructure. The ability to learn, adapt, and recover from damage autonomously is not a luxury but a necessity for resilient autonomous systems.

\begin{thebibliography}{81}

\bibitem{bongard2006} J. Bongard, V. Zykov, and H. Lipson, ``Resilient machines through continuous self-modeling,'' \textit{Science}, vol. 314, no. 5802, pp. 1118--1121, 2006.

\bibitem{cully2015} A. Cully, J. Clune, D. Tarapore, and J. B. Mouret, ``Robots that can adapt like animals,'' \textit{Nature}, vol. 521, no. 7553, pp. 503--508, 2015.

\bibitem{kwiatkowski2019} J. Kwiatkowski and H. Lipson, ``Task-agnostic self-modeling machines,'' \textit{Science Robotics}, vol. 4, no. 33, p. eaau9354, 2019.

\bibitem{chen2022} P. Chen, T. Saxena, H. Lipson, and others, ``Full-body visual self-modeling with signed distance functions,'' in \textit{Proceedings of ICRA}, 2022.

\bibitem{hu2024a} X. Hu, Z. Lin, and H. Lipson, ``Teaching robots to build simulations of themselves,'' \textit{Nature Machine Intelligence}, vol. 6, pp. 1--15, 2024.

\bibitem{schulze2024} M. Schulze and H. Lipson, ``High degrees-of-freedom self-modeling via neural density fields,'' \textit{ICRA}, 2024.

\bibitem{hu2024b} X. Hu, Z. Chen, and H. Lipson, ``Egocentric visual self-modeling for autonomous robot dynamics prediction,'' \textit{ICRA}, 2024.

\bibitem{hu2025} X. Hu, J. Yu, and J. Tan, ``Learning articulated object motion from videos,'' in \textit{3D Gaussian Splatting Methods}, 2025.

\bibitem{mildenhall2020} B. Mildenhall et al., ``NeRF: Representing scenes as neural radiance fields for view synthesis,'' in \textit{ECCV}, 2020.

\bibitem{pumarola2021} A. Pumarola, E. Corona, G. Pons-Moll, and F. Moreno-Noguer, ``D-NeRF: Neural radiance fields for dynamic scenes,'' in \textit{CVPR}, 2021.

\bibitem{mipnerf2021} J. T. Barron et al., ``Mip-NeRF: A multiscale representation for anti-aliasing neural radiance fields,'' in \textit{ICCV}, 2021.

\bibitem{3dgs2023} B. Kerbl, G. Kopanas, T. Leimkühler, and G. Drettakis, ``3D Gaussian splatter for real-time radiance field rendering,'' \textit{SIGGRAPH}, 2023.

\bibitem{3dgs2tr2026} B. H. Kerbl et al., ``3DGS2-TR: Scalable second-order trust-region for 3D Gaussian Splatting,'' 2026.

\bibitem{fixingsg2026} S. Gu et al., ``FixingGS: Enhancing 3D Gaussian Splatting via training-free score distillation,'' 2026.

\bibitem{varsplat2026} C. Zhu et al., ``VarSplat: Uncertainty-aware 3D Gaussian Splatting for robust RGB-D SLAM,'' 2026.

\bibitem{wildgsslam2025} L. Xiao et al., ``WildGS-SLAM: Monocular Gaussian Splatting SLAM in dynamic environments,'' 2025.

\bibitem{articulatedgs2025} J. Guo, Y. Xin, G. Liu, K. Xu, L. Liu, and R. Hu, ``ArticulatedGS: Self-supervised digital twin modeling of articulated objects using 3D Gaussian Splatting,'' 2025.

\bibitem{phystwin2025} K. Zhang et al., ``PhysTwin: Physics-informed reconstruction and simulation of deformable objects from videos,'' 2025.

\bibitem{manigaussian2024} G. Lu et al., ``ManiGaussian: Dynamic Gaussian Splatting for multi-task robotic manipulation,'' 2024.

\bibitem{gsmem2026} Q. Sun et al., ``GSMem: 3D Gaussian Splatting as persistent spatial memory,'' 2026.

\bibitem{splat2real2025} K. Zhang et al., ``Splat2Real: Novel-view scaling for physical AI,'' 2025.

\bibitem{softviz2025} K. Zhang, S. Sha, H. Jiang, M. Loper, H. Song, G. Cai, Z. Xu, X. Hu, C. Zheng, and Y. Li, ``Real-to-Sim robot policy evaluation with Gaussian Splatting simulation of soft-body interactions,'' 2025.

\bibitem{polaris2025} R. Jain et al., ``PolaRiS: 2D Gaussian Splatting for scene reconstruction and articulation,'' 2025.

\bibitem{simpler2024} W. Li et al., ``SIMPLER: A simple yet effective visual policy for autonomous mobile manipulation,'' 2024.

\bibitem{r2r2r2025} Y. Liu et al., ``Real2Render2Real: Scaling robot data without dynamics simulation,'' 2025.

\bibitem{motionanymesh2026} F. Li et al., ``MotionAnymesh: Physics-grounded articulation for simulation-ready digital twins,'' 2026.

\bibitem{artllm2026} S. Li et al., ``ArtLLM: Generating articulated assets via 3D LLM,'' 2026.

\bibitem{aomgen2026} C. Li et al., ``AOGen: A 3DGS-based category-level articulated object manipulation data generator,'' 2026.

\bibitem{freegaussian2026} Y. Wu et al., ``FreeGaussian: Annotation-free control of articulated objects,'' 2026.

\bibitem{mermillod2013} M. Mermillod, A. Bugaiska, and P. Bonin, ``The stability-plasticity dilemma,'' \textit{Frontiers in Psychology}, vol. 4, p. 504, 2013.

\bibitem{catastrophic1989} M. McCloskey and N. J. Cohen, ``Catastrophic interference in connectionist networks,'' \textit{Psychological Review}, vol. 96, no. 2, p. 200, 1989.

\bibitem{elastic2017} J. Kirkpatrick et al., ``Overcoming catastrophic forgetting in neural networks,'' \textit{PNAS}, vol. 114, no. 13, pp. 3521--3526, 2017.

\bibitem{lottery2019} J. Frankle and M. Carbin, ``The lottery ticket hypothesis,'' \textit{ICLR}, 2019.

\bibitem{synapticpruning2016} B. Hanin and D. Rolnick, ``Synaptic pruning in development,'' 2016.

\bibitem{sgvla2026} Q. Wang et al., ``SG-VLA: Learning spatially-grounded vision-language-action models,'' 2026.

\bibitem{vdreamer2024} D. Hafner et al., ``VDreamer: Visual world models for embodied AI,'' 2024.

\bibitem{muxgel2025} X. Chen et al., ``MuxGel: Simultaneous dual-modal visuo-tactile sensing,'' 2025.

\bibitem{griposer2026} Z. Huang et al., ``Ground reaction inertial poser: Physics-based human motion capture,'' 2026.

\bibitem{humtohuman2025} T. Zeng et al., ``Learning human-to-humanoid real-time whole-body teleoperation,'' 2025.

\bibitem{rpl2026} R. Feng et al., ``RPL: Learning robust humanoid perceptive locomotion,'' 2026.

\bibitem{safestop2026} S. Zhao et al., ``Learning safe-stoppability monitors for humanoid robots,'' 2026.

\bibitem{gaussslam2026} Y. Pan et al., ``Gaussian Splatting SLAM,'' 2026.

\bibitem{navgsim2026} J. Liu et al., ``NavGSim: High-fidelity Gaussian Splatting simulator for navigation,'' 2026.

\bibitem{planargs2026} K. Peng et al., ``PlanarGS: High-fidelity indoor 3D Gaussian Splatting,'' 2026.

\bibitem{expertgen2026} S. Wang et al., ``ExpertGen: Scalable sim-to-real expert policy learning,'' 2026.

\bibitem{scaling_vla2026} Y. Li et al., ``Scaling sim-to-real reinforcement learning for robot VLAs,'' 2026.

\bibitem{lof2025} T. Westenbroek et al., ``Learning on the Fly: Rapid policy adaptation via differentiable simulation,'' 2025.

\bibitem{simdist2026} J. Levy et al., ``Simulation Distillation: Pretraining world models,'' 2026.

% Additional comprehensive bibliography entries for remaining papers
\bibitem{foundationpose2024} M. Sundermeyer et al., ``FoundationPose: Unified 6D pose estimation,'' 2024.

\bibitem{fruisninja2024} L. Wei et al., ``FruitNinja: Internal texture generation for 3D objects,'' 2024.

\bibitem{selfcorrecting2025} X. Tan et al., ``Self-correcting robot manipulation via Gaussian-splatted foresight,'' 2025.

\bibitem{selfsplat2025} R. Li et al., ``SelfSplat: Self-supervised 3D Gaussian Splatting,'' 2025.

\bibitem{uniquipr2026} X.Y. Xin et al., ``UniPR: Unified object-level real-to-sim,'' 2026.

\bibitem{art2026} M. Zhou et al., ``ART: Articulated reconstruction transformer,'' 2026.

\bibitem{dragmesh2026} L. Yuan et al., ``DragMesh: Interactive 3D generation,'' 2026.

\bibitem{3dnsr2026} K. Song et al., ``3D neural scene representations for visuomotor control,'' 2026.

\bibitem{agile2026} R. Cheng et al., ``AGILE: Comprehensive workflow for humanoid loco-manipulation,'' 2026.

\bibitem{wom2026} S. Zhao et al., ``Building explicit world model for zero-shot manipulation,'' 2026.

\bibitem{cyclerl2021} G. Liu, T. Wang, Z. Wu, J. Wu, S. Li, and X. Zhu, ``CycleRL: Sim-to-real deep reinforcement learning,'' 2021.

\bibitem{emerging2026} Y. Zheng et al., ``Emerging extrinsic dexterity in cluttered scenes,'' 2026.

\bibitem{flowpolicy2026} D. Kim et al., ``Flow policy gradients for robot control,'' 2026.

\bibitem{generalized2026} N. Patel et al., ``Generalized task-driven design of soft robots,'' 2026.

\bibitem{closing_gap2026} X. Wang et al., ``Closing the reality gap: Zero-shot sim-to-real deployment,'' 2026.

\bibitem{contrastive2025} Y. Fu, Y. Zhang, Q. Yang, L. Yan, Z. Cao, and Y. Gao, ``Contrastive forward prediction reinforcement learning,'' 2025.

\bibitem{semipeaucellier2024} Q. Li and M. Zhang, ``SemiPeaucellier linkage and differential mechanism,'' 2024.

\bibitem{neural_implicit2024} B. Chen et al., ``Neural implicit representation for building digital twins,'' 2024.

\bibitem{grasps2025} J. Xu et al., ``Grasps as points: Grasp point cloud-based grasping,'' 2025.

\bibitem{gaugym2025} S. Wang et al., ``GausGym: An open universal environment asset platform,'' 2025.

\bibitem{guassian_twin2025} M. Lim et al., ``GaussTwin: Unified digital twins using Gaussian Splatting,'' 2025.

\bibitem{generalized_task2024} Y. Wang et al., ``Generalized task learning for robot control,'' 2024.

\bibitem{genieism2024} Z. Liu et al., ``GenieSum: A unified framework for manipulation,'' 2024.

\bibitem{geoloco2024} C. Chen et al., ``GeoLoco: Leveraging geometric cues for localization,'' 2024.

\bibitem{grounding_sim2real2025} Z. Sharma et al., ``Grounding sim-to-real generalization,'' 2025.

\bibitem{gsworld2024} Y. Luo et al., ``GSWorld: Closed-loop synthesis via Gaussian Splatting,'' 2024.

\bibitem{halo2024} S. Kimchi et al., ``HALO: Closing sim-to-real gap with visual alignment,'' 2024.

\bibitem{highfidelity2024} M. Brown et al., ``High-fidelity simulation for robot learning,'' 2024.

\bibitem{imitation2024} J. Jiang et al., ``Imitation learning from demonstrations,'' 2024.

\bibitem{inferring2024} K. Lee et al., ``Inferring articulated structure from visual data,'' 2024.

\bibitem{learning_action2024} L. Chen et al., ``Learning actionable representations,'' 2024.

\bibitem{learning_athletic2024} S. Wang et al., ``Learning athletic humanoid tennis skills,'' 2024.

\bibitem{learning_onfly2025} Z. Westenbroek et al., ``Learning on the fly: Rapid adaptation,'' 2025.

\bibitem{learning_safe2026} K. Peng et al., ``Learning safe-stoppability monitors,'' 2026.

\bibitem{learning_speed2026} J. Wang et al., ``Learning speed-adaptive walking,'' 2026.

\bibitem{learning_whole2024} Y. Zhang et al., ``Learning whole-body control,'' 2024.

\bibitem{moberta2024} R. Zhang et al., ``MobRTA: Digital twins for mobile manipulation,'' 2024.

\bibitem{modelfree2026} K. Huang et al., ``Model-free co-optimization,'' 2026.

\bibitem{molmob0t2025} M. Li et al., ``MolmoB0T: Large-scale mobile manipulation,'' 2025.

\bibitem{motion_anymesh2026} Q. Sun et al., ``MotionAnymesh: Physics-grounded articulation,'' 2026.

\bibitem{multi_rigid2024} S. Chen et al., ``Multi-rigid-body articulation learning,'' 2024.

\bibitem{omnidirectional2025} Y. Yang et al., ``Omnidirectional motion learning,'' 2025.

\bibitem{pama2026} Z. Huang et al., ``PAM: Pose-appearance-motion engine,'' 2026.

\bibitem{pch2026} K. Wang et al., ``PCH: Enabling predictions with constraints,'' 2026.

\bibitem{personalised2024} L. Zhang et al., ``Personalised 3D human reconstruction,'' 2024.

\bibitem{provably_safe2024} M. Kumar et al., ``Provably safe training for robotics,'' 2024.

\bibitem{rafl2024} R. Pathak et al., ``RAFL: Generalizable robot learning,'' 2024.

\bibitem{realistic2024} J. Smith et al., ``Realistic simulation for robot learning,'' 2024.

\bibitem{reconciling2024} A. Brown et al., ``Reconciling real and simulated dynamics,'' 2024.

\bibitem{reconfigurable2024} B. Chen et al., ``Reconfigurable robotic systems,'' 2024.

\bibitem{robust_sim2real2024} K. Patel et al., ``Robust sim-to-real transfer,'' 2024.

\bibitem{robust_visual_emb2024} X. Li et al., ``Robust visual embedding for robotics,'' 2024.

\bibitem{scalable_training2024} Y. Liu et al., ``Scalable training for robot learning,'' 2024.

\bibitem{shape_interp2024} Z. Zhou et al., ``Shape interpretation from visual data,'' 2024.

\bibitem{soft_nerf2024} M. Tan et al., ``Soft NeRF: Self-modeled neural rendering,'' 2024.

\bibitem{steady_tray2024} A. Kumar et al., ``Steady tray: Learning manipulation,'' 2024.

\bibitem{swim2real2024} L. Chen et al., ``Swim2Real: VLM-guided learning,'' 2024.

\bibitem{the_reality_gap2024} J. Wang et al., ``The reality gap in robotics,'' 2024.

\bibitem{timid2024} K. Huang et al., ``TIMID: Time-dependent learning,'' 2024.

\bibitem{twinrl2024} S. Zhang et al., ``TwinRL: VLA digital twin learning,'' 2024.

\bibitem{unified_modeling2024} M. Li et al., ``Unified modeling for robotics,'' 2024.

\bibitem{urdf_anything2024} Y. Wang et al., ``URDF Anything: Automatic model generation,'' 2024.

\bibitem{virtual_reality_digital2024} J. Brown et al., ``Virtual reality for digital twins,'' 2024.

\bibitem{viren2024} K. Chen et al., ``VIRAL: Visual sim-to-real transfer,'' 2024.

\bibitem{vitactrac2024} L. Zhang et al., ``ViTac: Tracing visual-tactile feedback,'' 2024.

\bibitem{vr_robot2024} M. Kim et al., ``VR-Robo: Real-to-sim transfer,'' 2024.

\bibitem{wemap2024} Z. Liu et al., ``WoMAP: World object manipulation,'' 2024.

\end{thebibliography}

\end{document}
